\section{Experimental design and evidence}
The four experiments count a result only after its correctness invariant is checked. E1 and E2 use 20 independently generated RSA key-set blocks; the three E2 messages within each block are paired observations, not independent key sets. Cryptographic randomness was not seeded; the recorded schedule seed controls workload order only.

E1's four messages are \texttt{B} repeated 1, 96 or 180 times, and 95 zero bytes followed by \texttt{0x02}; the two 96-byte samples test the distinction between byte count and numeric value. Each of the 20 blocks evaluates all seven subsets for every sample, including every pair. E2 reuses the block's keys in both schemes, using the 1-, 96- and 180-byte \texttt{B} samples. Textbook and OAEP ciphertexts are each submitted to the same attacker implementation. OAEP encryption is additionally repeated under one public key and every intended recipient's decryption is checked. This design can attribute a within-pair difference to the scheme change more clearly than a comparison using independently generated keys, while retaining the limitation that 60 pairs are clustered within 20 blocks.

E3 prepares a five-recipient public-key set before routine timing, uses its first 1, 3 or 5 members and encrypts 10\,KiB, 100\,KiB and 1\,MiB synthetic files with a \emph{fresh} AES key/nonce on every repetition. Two warmups per configuration are excluded; 20 measured repetitions remain for each of nine configurations. The monotonic \texttt{perf\_counter\_ns} interval \texttt{encryption\_ns} spans key generation for AES, RSA-OAEP key wrapping, AAD creation and AES-GCM content encryption. \texttt{serialization\_ns} spans canonical JSON/Base64 output; \texttt{total\_encryption\_ns} is their sum. RSA recipient-key generation is measured separately in 60 records, so it is not silently included in routine encryption. \texttt{decryption\_ns} measures the first valid recipient only. Every decrypted byte string is compared with the input, and every serialized size is measured from actual package bytes. Reported summaries use all individual timings, with means, medians, sample standard deviations and quartiles; no outlier was removed.

The official manifest records Windows~11, Python~3.12.10, \texttt{cryptography}~48.0.1, an Intel64 Family~6 Model~191 CPU with 20 logical CPUs, run ID \texttt{official-20261106T083221Z-3bbf57cc}, and source commit \texttt{46c4558}. Hardware-specific timing must not be generalized to cloud deployments. E4 checks one valid package and one mutation for each of ten rejection conditions; it records both the observed category and whether any plaintext escaped. The archived \href{https://github.com/Dung092005/Crypto-project/tree/main/docs/evidence/official-20261106T083221Z-3bbf57cc}{official evidence} and \href{https://github.com/Dung092005/Crypto-project/blob/main/docs/REPRODUCIBILITY.md}{reproduction procedure} support read-back validation without repeating RSA key generation or benchmarks.
\FloatBarrier
